%% ma: L11-B09-0001
%% lop: 11
%% chuong: 3
%% bai: 9
%% dang: 11.9.1
%% loai: TN4
%% muc: TH
%% dap_an: "C"
%% nguon: "(NBV)-ĐỀ SỐ 1-MỖI NGÀY 1 ĐỀ THI 2025.pdf (D01, MỖI NGÀY 1 ĐỀ - Đề số 1), Phần I Câu 4"
%% kien_thuc: "Lớp 11 Bài 9 (số trung bình của mẫu số liệu ghép nhóm)"
%% ghi_chu: "Mức TH (tính hai số trung bình theo công thức quen thuộc rồi so sánh); có thể nhận xét nhanh vì M_2 là M_1 tịnh tiến 6 đơn vị"
%% trang_thai: chinh-thuc
\begin{ex}
Cho hai mẫu số liệu ghép nhóm $M_1$, $M_2$ có bảng tần số ghép nhóm như sau:
\begin{center}
\renewcommand{\arraystretch}{1.2}
\begin{tabular}{|c|c|c|c|c|c|}
\hline
$M_1$: Nhóm & $[12;14)$ & $[14;16)$ & $[16;18)$ & $[18;20)$ & $[20;22)$ \\
\hline
Tần số & $2$ & $4$ & $8$ & $6$ & $4$ \\
\hline
\end{tabular}

\medskip
\begin{tabular}{|c|c|c|c|c|c|}
\hline
$M_2$: Nhóm & $[6;8)$ & $[8;10)$ & $[10;12)$ & $[12;14)$ & $[14;16)$ \\
\hline
Tần số & $2$ & $4$ & $8$ & $6$ & $4$ \\
\hline
\end{tabular}
\end{center}
Gọi $\overline{X}_1$ và $\overline{X}_2$ lần lượt là số trung bình của mẫu số liệu ghép nhóm $M_1,M_2$. Phát biểu nào sau đây là đúng?
\choice{$\overline{X}_1=\overline{X}_2$}{$\overline{X}_1=2\overline{X}_2$}{\True $\overline{X}_1=\overline{X}_2+6$}{$2\overline{X}_1=\overline{X}_2$}
\loigiai{Cả hai mẫu đều có $n=24$. Giá trị đại diện của $M_1$: $13;15;17;19;21$, nên $\overline{X}_1=\dfrac{13\cdot2+15\cdot4+17\cdot8+19\cdot6+21\cdot4}{24}=\dfrac{420}{24}=17{,}5$.
Giá trị đại diện của $M_2$: $7;9;11;13;15$, nên $\overline{X}_2=\dfrac{276}{24}=11{,}5$.
Vậy $\overline{X}_1=\overline{X}_2+6$.}
\end{ex}
